\scanpage{56}
% archive sheet number: 55
\mspage{2 - 4 -}\hfill n\textsuperscript{o} 392

$1\handwritten{\leq}i\handwritten{\leq}n$. Soit $x_{\unsure{0|o}}$ un élément de\unsure{,} $K$ tel que $v_i(x_0)=q_i$ pour $1\handwritten{\leq}i\handwritten{\leq}n$, (chap. VI, § 7, n\textsuperscript{o} 2, cor. 1 du th. 1). Alors $\underset{=}{a}$ est l'ensemble des $x\handwritten{\in}K$ tels que $v_i(xx_0^{-1})\handwritten{\geq}0$ pour $1\handwritten{\leq}i\handwritten{\leq}n$. Ainsi $\underset{=}{a}=Ax_0$.

Si $A$ est un anneau de Dedekind, on a vu, dans la démonstration du th. 1, que le groupe $D(A)$ des diviseurs de $A$ s'identifie au groupe $I(A)$ des idéaux fractionnaires $\underset{=}{a}\handwritten{\ne}(0)$, (comme $A$ est noethérien, tout idéal fractionnaire non nul est de type fini). Le groupe des \uline{classes de diviseurs} de $A$, (§ 1, n\textsuperscript{o} 4), s'identifie alors au groupe des \uline{classes d'idéaux} de $A$, (chap. II, § 5, n\textsuperscript{o} 7).

3. \uline{Décomposition des idéaux en produits d'idéaux premiers}.

Soient $A$ un anneau de Dedekind, $I(A)$ le groupe multiplicatif ordonné des idéaux fractionnaires non nuls de $A$, et $D(A)$ le groupe des diviseurs de $A$. L'isomorphisme $\underset{=}{a}\longrightarrow\operatorname{div}\underset{=}{a}$ de $I(A)$ sur $D(A)$ fait correspondre les diviseurs extrémaux aux idéaux premiers non nuls de $A$, (§ 1, n\textsuperscript{o} 6, th. 3), donc le groupe multiplicatif $I(A)$ admet pour base l'ensemble des idéaux premiers non nuls de $A$, (§ 1, n\textsuperscript{o} 3, th. 2). Autrement dit, \uline{tout idéal fractionnaire non nul} $\underset{=}{a}$ \uline{de} $A$ \uline{admet une décomposition et une seule de la forme} :
\[
\underset{=}{a}=\handwritten{\prod}\unsure{\underset{=}{p}|\underline{p}|p}^{n(\underset{=}{p})}\tag*{(1)}
\]
où le produit est étendu aux idéaux premiers non nuls de $A$, les exposants $n(\underset{=}{p})$ étant nuls à l'exception d'un nombre fini d'entre eux. De plus $\underset{=}{a}$ est entier si et seulement si les $n(\underset{=}{p})$ sont tous positifs. On dit que la relation (1) est la \uline{décomposition de} $\underset{=}{a}$ \uline{en facteurs premiers}. En particulier, si $\underset{=}{a}$ est un idéal principal $Ax$, on a, pour tout $\underset{=}{p}$, $n(\underset{=}{p})=v_{\underset{=}{p}}(x)$, où $v_{\underset{=}{p}}$ désigne la valuation essentielle correspondant à $\underset{=}{p}$ ; ceci résulte en effet de la formule (4) du § 1, n\textsuperscript{o} 4. Soient
\[
\underset{=}{a}=\handwritten{\prod}\underset{=}{p}^{m(\underset{=}{p})}\quad,\qquad \underset{=}{b}=\handwritten{\prod}\underset{=}{p}^{n(\underset{=}{p})}
\]
deux idéaux fractionnaires non nuls de $A$. On a alors :
